\par\begingroup
\setlength{\lineskip}{4pt}\setlength{\lineskiplimit}{4pt}
\Needspace{8\baselineskip}\bigskip
\pdfbookmark[3]{Exposition, 18}{pallas-II-article18}
\begin{center}18.\end{center}
\noindent\hspace*{\parindent}Let \(\mathrm{f}t\) be a function of \(t\) whose numerical values are given for \(t=T,\ t=T+\theta,\ t=T+2\theta,\ t=T+3\theta\) etc. Let us denote by \(\mathrm{f}'t\) the difference
\[
\mathrm{f}(t+\tfrac12\theta)-\mathrm{f}(t-\tfrac12\theta),
\]
and similarly
\[
\begin{aligned}
\text{by }\mathrm{f}''t&\text{ the difference }\mathrm{f}'(t+\tfrac12\theta)-\mathrm{f}'(t-\tfrac12\theta)\\
\mathrm{f}'''t&\text{ the difference }\mathrm{f}''(t+\tfrac12\theta)-\mathrm{f}''(t-\tfrac12\theta)\\
\mathrm{f}^{\mathrm{IV}}t&\text{ the difference }\mathrm{f}'''(t+\tfrac12\theta)-\mathrm{f}'''(t-\tfrac12\theta)
\end{aligned}
\]
and so on.
\par Forming the series of differences gives us the values of \SourceLeaf{468}\(\mathrm{f}'(T+\tfrac12\theta),\ \mathrm{f}'(T+\tfrac32\theta),\ \mathrm{f}'(T+\tfrac52\theta)\) etc., \(\mathrm{f}''(T+\theta),\ \mathrm{f}''(T+2\theta),\ \mathrm{f}''(T+3\theta)\) etc., \(\mathrm{f}'''(T+\tfrac32\theta),\ \mathrm{f}'''(T+\tfrac52\theta)\) etc. etc. Let us suppose further that
\[
\begin{gathered}
\begin{aligned}
\mathrm{f}t&=\mathrm{f}^{(-1)}(t+\tfrac12\theta)-\mathrm{f}^{(-1)}(t-\tfrac12\theta)\\
\mathrm{f}^{(-1)}t&=\mathrm{f}^{(-2)}(t+\tfrac12\theta)-\mathrm{f}^{(-2)}(t-\tfrac12\theta)
\end{aligned}\\
\text{etc.}
\end{gathered}
\]
An arbitrary value may be assigned to \(\mathrm{f}^{(-1)}(T-\tfrac12\theta)\), and the successive addition of the values of \(\mathrm{f}t\) will give the values of \(\mathrm{f}^{(-1)}(T+\tfrac12\theta),\ \mathrm{f}^{(-1)}(T+\tfrac32\theta),\ \mathrm{f}^{(-1)}(T+\tfrac52\theta)\) etc. Similarly, an arbitrary value will be assigned to \(\mathrm{f}^{(-2)}(T-\theta)\), and the successive addition of the values of \(\mathrm{f}^{(-1)}t\) will give the values of \(\mathrm{f}^{(-2)}T,\ \mathrm{f}^{(-2)}(T+\theta),\ \mathrm{f}^{(-2)}(T+2\theta)\) etc., and so on.
\par With these assumptions we have
\[
\begin{aligned}
\mathrm{I.}\quad&\int\mathrm{f}t.\,\mathrm{d}t=\alpha+\theta\bigl\{\mathrm{f}^{(-1)}t+\tfrac1{24}\mathrm{f}'t\\
&\qquad-\tfrac{17}{5760}\mathrm{f}'''t+\tfrac{367}{967680}\mathrm{f}^{\mathrm{V}}t\\
&\qquad-\tfrac{27859}{464486400}\mathrm{f}^{\mathrm{VII}}t+\tfrac{1295803}{122624409600}\mathrm{f}^{\mathrm{IX}}t-\text{etc.}\bigr\}
\end{aligned}
\]
\[
\begin{aligned}
\mathrm{II.}\quad&\int\!\int\mathrm{f}t.\,\mathrm{d}t^2=\alpha t+\beta+\theta\theta\bigl\{\mathrm{f}^{(-2)}t+\tfrac1{12}\mathrm{f}t\\
&\qquad-\tfrac1{240}\mathrm{f}''t+\tfrac{31}{60480}\mathrm{f}^{\mathrm{IV}}t\\
&\qquad-\tfrac{289}{3628800}\mathrm{f}^{\mathrm{VI}}t+\tfrac{317}{22809600}\mathrm{f}^{\mathrm{VIII}}t-\text{etc.}\bigr\},
\end{aligned}
\]
where \(\alpha\) and \(\beta\) are constant quantities.
\par The coefficients of the first series are obtained by dividing unity by
\[
\begin{aligned}
&1-\tfrac1{24}xx+\tfrac3{640}x^4-\tfrac5{7168}x^6+\text{etc.}\\
&\quad=1-\tfrac18\cdot\tfrac13xx+\dfrac{1.3}{8.16}\cdot\tfrac15\cdot x^4\\
&\qquad\quad-\dfrac{1.3.5}{8.16.24}\cdot\tfrac17\cdot x^6+\dfrac{1.3.5.7}{8.16.24.32}\cdot\tfrac19\cdot x^8-\text{etc.}\\
&\quad=\dfrac{\log(\BareRoot{(1+\tfrac14xx)}+\tfrac12x)}{\tfrac12x}\mathpunct{\cdot}
\end{aligned}
\]
Squaring the quotient gives the coefficients of the second series. Higher powers of the quotient would likewise give the coefficients of analogous series for expressing the integrals \(\int^3\mathrm{f}t.\,\mathrm{d}t^3\) etc., which we do not need for our investigation.
\par Here are two further series that can be used:
\[
\begin{aligned}
\mathrm{III.}\quad&\int\mathrm{f}t.\,\mathrm{d}t=\alpha+\tfrac12\theta\bigl\{\\
&\quad(\mathrm{f}^{(-1)}(t-\tfrac12\theta)+\mathrm{f}^{(-1)}(t+\tfrac12\theta))\\
&\quad-\tfrac1{12}(\mathrm{f}'(t-\tfrac12\theta)+\mathrm{f}'(t+\tfrac12\theta))\\
&\quad+\tfrac{11}{720}(\mathrm{f}'''(t-\tfrac12\theta)+\mathrm{f}'''(t+\tfrac12\theta))\\
&\quad-\tfrac{191}{60480}(\mathrm{f}^{\mathrm{V}}(t-\tfrac12\theta)+\mathrm{f}^{\mathrm{V}}(t+\tfrac12\theta))+\text{etc.}\bigr\}
\end{aligned}
\]
\[
\begin{aligned}
\mathrm{IV.}\quad&\int\!\int\mathrm{f}t.\,\mathrm{d}t^2=\alpha t+\beta+\tfrac12\theta\theta\bigl\{\\
&\quad(\mathrm{f}^{(-2)}(t-\tfrac12\theta)+\mathrm{f}^{(-2)}(t+\tfrac12\theta))\\
&\quad-\tfrac1{24}(\mathrm{f}(t-\tfrac12\theta)+\mathrm{f}(t+\tfrac12\theta))\\
&\quad+\tfrac{17}{1920}(\mathrm{f}''(t-\tfrac12\theta)+\mathrm{f}''(t+\tfrac12\theta))\\
&\quad-\tfrac{367}{193536}(\mathrm{f}^{\mathrm{IV}}(t-\tfrac12\theta)+\mathrm{f}^{\mathrm{IV}}(t+\tfrac12\theta))+\text{etc.}\bigr\}.
\end{aligned}
\]
\SourceLeaf{469}The coefficients of the third series are obtained by dividing \(1-\tfrac18xx+\dfrac{1.3}{8.16}x^4-\dfrac{1.2.5}{8.16.24}x^6+\text{etc.}\) by \(1-\tfrac18\cdot\tfrac13xx+\dfrac{1.3}{8.16}\cdot\tfrac15x^4-\dfrac{1.3.5}{8.16.24}\cdot\tfrac17x^6+\text{etc.}\) or by expanding the following in a series
\[
\dfrac{\tfrac12x}{\BareRoot{(1+\tfrac14xx)}.\log(\BareRoot{(1+\tfrac14xx)}+\tfrac12x)}\mathpunct{\cdot}
\]
Similarly, the coefficients of the fourth series arise from the expansion of
\[
\dfrac{\tfrac14xx}{\BareRoot{(1+\tfrac14xx)}.\{\log(\BareRoot{(1+\tfrac14xx)}+\tfrac12x)\}^2}\mathpunct{\cdot}
\]
The coefficients of the fourth series are obtained still more easily by multiplying those of the first by \(1,\ -1,\ -3,\ -5,\ -7\) etc. Integrals of higher order could also be expressed in an analogous way.
\par We omit the proofs of these theorems here, since they can easily be deduced from the theory of generating functions recently presented by an illustrious geometer.
\par\endgroup
